\scanpage{82}
\mspage{5}
% archive foliation: 81
\[
(30)\qquad T_{E/S,0}^{\otimes2}\simeq T_{\Sigma_{\underline{J}},t}
\]
où $t$ est la section de $\Sigma_J^*=\Sigma_J\setminus D_J$ sur $S$, l'image de la section nulle de $E$. Ceci posé, on définit un \uline{épinglage de Legendre} de $E/S$ comme étant formé d'un iso
\[
(31)\quad\deleted{\unsure{T}}\underbrace{\Delta_{E/S}}_{T_{E/S,0}}\ \deleted{\unsure{\alpha}}\overset{\kappa}{\underset{\sim}{\longrightarrow}}\ \deleted{\unsure{T}}\underline{O}_{\deleted{\unsure{\Sigma_J}}}(1)_{t}^{\unsure{\sharp}}\simeq
\]
tel que \deleted{il existe un} l'isom\inserted{$\kappa^{\unsure{\otimes2|2}}$} déduit par passage au carré tensoriel s'identifie, via (30) et (27), à l'isom
\[
(32)\qquad T_{\Sigma_J,t}\simeq T_{\Sigma_J,t}\wedge_{\{\pm1\}}(\unsure{\boldsymbol{\mu}|\mu}_4^*\wedge\underline{\omega}_{\underline{J}})
\]
déduit d'une section
\[
(33)\quad\left\{\begin{array}{l}
\alpha\in\Gamma(S,\unsure{\boldsymbol{\mu}|\mu}_4^*\wedge\underline{\omega}_{\underline{J}})\qquad\text{i.e. d'un \uline{iso}}\\[5pt]
\alpha:\unsure{\boldsymbol{\mu}|\mu}_4^*\simeq\underline{\omega}_{\underline{J}}
\end{array}\right.
\]
(qui est déterminé de façon unique en termes
